\section{Introduction}
Recent advancements in Large Language Models (LLMs) have significantly impacted various fields thanks to their exceptional ability to integrate knowledge and follow instructions~\cite{achiam2023gpt,touvron2023llama}
Leveraging techniques such as chain-of-thought reasoning~\cite{wei2022chain}, Retrieval-Augmented Generation (RAG)~\cite{lewis2020retrieval}, and multi-agent collaboration frameworks~\cite{park2023generative}, LLMs show great potential in creative tasks, including software development~\cite{hong2023metagpt}, neural architecture search~\cite{zheng2023can}, hyperparameter optimization~\cite{zhang2023using,yang2024autommlab}, and scientific discovery~\cite{lu2024aiscientist}.
\begin{figure}
    \centering
    \includegraphics[width=0.49\textwidth]{figures/teaser.pdf}
    \caption{Traditional Neural Architecture Search (NAS) approaches aim to discover high-performing neural architectures within an expert-defined search space. In contrast, Neural Architecture Design (NAD) is not constrained by predefined search spaces, allowing for the exploration of entirely new architectures.
    }
    \label{fig:teaser}
\end{figure}
Designing neural network architectures that meet specific performance criteria (\eg, high accuracy or low latency) remains a significant challenge in deep learning. This process often involves extensive manual efforts and costly trial-and-error cycles. To address these challenges, existing studies have leveraged Neural Architecture Search (NAS) to automate the discovery of high-performing architectures, achieving notable results across diverse domains~\cite{zoph2016neural,liu2018darts}. However, traditional NAS methods typically rely on expert-defined search spaces, limiting their capacity to uncover truly novel architectures.

Motivated by advancements in large language models (LLMs), recent research has explored a new direction in Neural Architecture Design (NAD) driven by LLMs, moving beyond conventional predetermined search spaces. Unlike traditional NAS approaches, NAD fosters innovation by enabling the exploration of entirely new architectures within an open architecture space. However, this endeavor presents several challenges.
First, many existing studies rely on evolutionary computation or reinforcement learning for architecture modifications, which often result in inefficient, undirected exploration, leading to training and evaluation of an overwhelming number of neural network architectures.  
Second, existing LLM-based methods perform NAD tasks independently, without benefiting from immediate feedback or past experiences, leading to repeated mistakes and unnecessary trail-and-error processes.
Additionally, code-generating LLMs often struggle to maintain focus on neural architecture design, becoming sidetracked by non-essential aspects like text style and code syntax.
Such distractions can significantly diminish the efficiency of the automatic exploration.

To tackle these challenges, we propose NADER (Neural Architecture Design via multi-Agent collaboRation), an LLM-driven NAD method. NADER begins with a well-established network, such as ResNet~\cite{he2016deep}, and seeks to enhance it through iterative modifications. 
The framework comprises a research team and a development team, each consisting of multiple collaborating agents. This multi-agent approach facilitates more efficient and directed exploration. 
The research team is responsible for proposing improvement proposals based on insights drawn from recent academic literature and model optimization trees.
Each proposal contains a model to be improved and a modification suggestion.
Subsequently, the development team implements these proposals and evaluate the performance of the modified network.
The performance is then relayed back to the research team, informing subsequent proposals and fostering a continuous cycle of improvement.

``Error is often the correct guide." We introduce the Reflector, which effectively incorporates immediate feedback from the external environment and leverages historical experiences.
Reflector tracks the behavior of Modifier and reflect on it to gain experience.
Given the proposed modification suggestions, the Reflector queries the experience pool to retrieve relevant experience, providing the Modifier with contextual information for implementation.
This reduces the likelihood of repeating past mistakes and increases the success rate of a good modifications.

Furthermore, to help the development team concentrate on network structure, we propose a graph-based network representation. Instead of generating executable code directly, NADER represents neural architectures as single directed acyclic graphs (DAGs), where nodes represent operations (\eg, convolution, normalization) and edges denote information flow. This graph representation allows NADER to prioritize high-level structural decisions, avoiding distractions from low-level code implementation details.

We demonstrate the effectiveness of NADER in discovering high-performing neural architectures that extend beyond predetermined search spaces. Through extensive experiments, we showcase the advantages of NADER over state-of-the-art NAS and NAD methods. Our contributions are summarized as follows:
\begin{itemize}
    \item We introduce NADER, a novel LLM-driven framework for neural architecture design that utilizes a multi-agent collaboration strategy to automatically explore novel neural architectures beyond predetermined search spaces. 
    \item  We introduce Reflector, which effectively incorporates immediate feedback from the external environment and experiences reflected from historical trajectories, thereby reducing repeated mistakes and enhancing efficiency.
    \item We introduce a graph-based representation for neural network architectures, facilitating a focus on high-level structural decisions while minimizing distractions from low-level coding details.
\end{itemize}




